% =====================================================================
%  INTEGRA Sono | Piloto de saúde do sono via WhatsApp (NR-1 + LGPD)
%  Compilação: pdflatex -> biber -> pdflatex -> pdflatex
%  (ou: latexmk -pdf Apresentacao.tex)
% =====================================================================
\documentclass[aspectratio=169,10pt]{beamer}

% ---------- Tema e fontes ----------
\usetheme[progressbar=frametitle,numbering=fraction,block=fill,
          sectionpage=progressbar,titleformat section=regular]{metropolis}
\usepackage[T1]{fontenc}
\usepackage[utf8]{inputenc}
\usepackage[brazilian]{babel}
\usepackage[sfdefault,lining]{FiraSans}
\usepackage[scale=0.85]{FiraMono}
\usepackage{csquotes}

% ---------- Tabelas, ícones, figuras ----------
\usepackage{booktabs,tabularx,array,multirow,makecell}
\usepackage{adjustbox}
\usepackage{fontawesome5}
\usepackage{tikz}
\usetikzlibrary{arrows.meta,positioning,shapes.geometric,shapes.misc,
                shapes.multipart,calc,fit,backgrounds,patterns,shadows.blur}
\usepackage{pgfplots}
\pgfplotsset{compat=1.18}
\usepackage{pgfgantt}
\usepackage[most]{tcolorbox}

% ---------- Bibliografia ----------
\usepackage[backend=biber,style=authoryear-comp,maxcitenames=2,
            uniquelist=false,uniquename=false,dashed=false,
            giveninits=true,doi=true,url=true,isbn=false]{biblatex}
\usepackage{xurl}
\addbibresource{Apresentacao.bib}
% Separa citações múltiplas com vírgula em vez de ponto e vírgula
\renewcommand*{\multicitedelim}{\addcomma\space}
\renewcommand*{\bibfont}{\tiny}
\setlength{\bibitemsep}{0.3\baselineskip}
\setbeamertemplate{bibliography item}{}
\setbeamerfont{bibliography item}{size=\tiny}
\setbeamerfont{bibliography entry author}{size=\tiny}
\setbeamerfont{bibliography entry title}{size=\tiny}
\setbeamerfont{bibliography entry location}{size=\tiny}
\setbeamerfont{bibliography entry note}{size=\tiny}
\setbeamercolor{bibliography entry author}{fg=grafite}
\setbeamercolor{bibliography entry title}{fg=noite}
\setbeamercolor{bibliography entry location}{fg=mudo}
\setbeamercolor{bibliography entry note}{fg=mudo}
% URL só quando não há DOI
\AtEveryBibitem{\iffieldundef{doi}{}{\clearfield{url}\clearfield{urldate}}}

% ---------- Paleta ----------
\definecolor{noite}{HTML}{1D2B4A}     % azul noturno (cor principal)
\definecolor{lavanda}{HTML}{6C5BA7}   % sono / recuperação
\definecolor{menta}{HTML}{2A9D8F}     % positivo / saúde
\definecolor{ambar}{HTML}{E9A23B}     % atenção
\definecolor{coral}{HTML}{C8553D}     % alerta
\definecolor{nevoa}{HTML}{F3F4F8}     % fundos claros
\definecolor{grafite}{HTML}{3A3F4B}   % texto
\definecolor{mudo}{HTML}{7A8091}      % texto secundário
\definecolor{wa}{HTML}{075E54}        % cabeçalho WhatsApp
\definecolor{wabg}{HTML}{ECE5DD}      % fundo do chat
\definecolor{wauser}{HTML}{DCF8C6}    % balão do usuário

\setbeamercolor{normal text}{fg=grafite,bg=white}
\setbeamercolor{frametitle}{fg=white,bg=noite}
\setbeamercolor{progress bar}{fg=ambar,bg=noite!40}
\setbeamercolor{title separator}{fg=ambar}
\setbeamercolor{alerted text}{fg=coral}
\setbeamercolor{example text}{fg=menta}
\setbeamercolor{block title}{fg=white,bg=noite}
\setbeamercolor{block body}{fg=grafite,bg=nevoa}
\setbeamercolor{block title alerted}{fg=white,bg=coral}
\setbeamercolor{block title example}{fg=white,bg=menta}
\setbeamercolor{footnote}{fg=mudo}
\setbeamerfont{footnote}{size=\tiny}
\setbeamertemplate{itemize item}{\color{lavanda}\small\faCircle[regular]}
\setbeamertemplate{itemize subitem}{\color{mudo}\tiny\faChevronRight}

% ---------- Macros ----------
% Linha de fonte no rodapé do slide
\newcommand{\fonte}[1]{\par\vfill{\tiny\color{mudo}\faBookOpen\ #1\par}}
% Ilustração externa: insere a imagem se existir; senão desenha um
% espaço reservado com o nome do arquivo a ser baixado.
% Uso: \imagem{largura}{altura}{img/arquivo.png}{descrição}
\newcommand{\imagem}[4]{%
  \IfFileExists{#3}%
    {\includegraphics[width=#1,height=#2,keepaspectratio]{#3}}%
    {\begin{tikzpicture}
       \node[draw=mudo,dashed,rounded corners=4pt,fill=nevoa,
             minimum width=#1,minimum height=#2,text width=#1-6mm,
             align=center,font=\tiny\color{mudo}]
         {{\Large\faImage}\\[2pt]\texttt{\detokenize{#3}}\\[1pt]\textit{#4}};
     \end{tikzpicture}}%
}
% Cartão de indicador (número grande + legenda)
\newcommand{\kpi}[4]{% cor, ícone, número, legenda
  \begin{tcolorbox}[enhanced,colback=white,colframe=#1,boxrule=0pt,
      leftrule=3pt,arc=2pt,left=6pt,right=4pt,top=4pt,bottom=4pt,
      height=3.35cm,valign=top]
    {\color{#1}\large #2}\par\smallskip
    {\color{#1}\LARGE\bfseries #3}\par\smallskip
    {\scriptsize\raggedright #4\par}
  \end{tcolorbox}}
% Marco da linha do tempo: x, cor, ícone, data, descrição, above/below
\newcommand{\marco}[6]{%
  \node[circle,fill=#2,draw=white,line width=1.2pt,minimum size=5mm,
        inner sep=0pt,font=\tiny\color{white}] (m) at (#1,0) {#3};
  \node[text width=3.2cm,align=center,font=\scriptsize,execute at begin node={\hyphenpenalty=10000},#6=3mm of m]
    {\textbf{\color{#2}#4}\\ #5};}
\newcommand{\nao}{\textcolor{coral}{\faTimesCircle}}
\newcommand{\tagc}[2]{\tikz[baseline=(t.base)]\node[fill=#1,text=white,
  rounded corners=2pt,inner sep=2pt,font=\tiny\bfseries](t){#2};}

\newcolumntype{L}{>{\raggedright\arraybackslash}X}
\newcolumntype{C}{>{\centering\arraybackslash}X}
\renewcommand{\arraystretch}{1.25}

% ---------- Estilos TikZ ----------
\tikzset{
  caixa/.style={draw=#1,fill=#1!8,rounded corners=3pt,line width=0.7pt,
     align=center,font=\scriptsize,inner sep=4pt},
  caixa/.default=noite,
  seta/.style={-{Stealth[length=2.2mm]},line width=0.8pt,color=#1},
  seta/.default=grafite,
  grupo/.style={draw=#1,dashed,rounded corners=6pt,inner sep=7pt,
     line width=0.6pt},
  rot/.style={font=\tiny\bfseries\color{#1},fill=white,inner sep=1.5pt},
  etapa/.style={circle,fill=#1,text=white,minimum size=17mm,
     font=\scriptsize\bfseries,align=center,inner sep=1pt},
  passo/.style={rounded corners=8pt,fill=#1,text=white,minimum width=2.1cm,
     minimum height=7mm,font=\scriptsize\bfseries},
  ev/.style={caixa=#1,text width=3.0cm,minimum height=1.55cm},
  pas/.style={caixa=#1,text width=2.35cm,minimum height=2.0cm},
  cx/.style={caixa=#1,text width=2.3cm,minimum height=0.85cm},
  db/.style={cylinder,shape border rotate=90,aspect=0.18,draw=#1,fill=#1!10,
     minimum width=2.5cm,minimum height=0.95cm,align=center,font=\tiny},
  cab/.style={caixa=#1,text width=2.1cm,minimum height=0.8cm},
  msg/.style={-{Stealth[length=2mm]},line width=0.7pt,color=#1},
  msg/.default=grafite,
  ent/.style={rectangle split,rectangle split parts=2,draw=#1,
     rectangle split part fill={#1,white},rounded corners=2pt,
     font=\tiny,text width=2.9cm,align=left,inner sep=3pt},
  st/.style={caixa=#1,text width=2.0cm,minimum height=1.75cm},
  dbn/.style={cylinder,shape border rotate=90,aspect=0.2,draw=#1,fill=#1!10,
     text width=2.0cm,minimum height=1.1cm,align=center,font=\scriptsize},
  st2/.style={caixa=#1,text width=2.3cm,minimum height=1.7cm},
}

% ---------- Metadados ----------
\title{INTEGRA Sono}
\subtitle{Piloto de saúde do sono no trabalho via WhatsApp com painel de
  gestão alinhado a NR-1}
\author{Programa INTEGRA \textperiodcentered\ Equipe MorpheusIA}
\institute{Proposta técnica para discussão}
\date{Outubro de 2026}
\titlegraphic{\begin{tikzpicture}[remember picture,overlay]\node[anchor=south east,inner sep=0pt] at ([xshift=-0.9cm,yshift=0.8cm]current page.south east){\imagem{5.2cm}{3.4cm}{img/capa_trabalhador_amanhecer.png}{Pessoa acordando descansada ao amanhecer (Foto: Miriam Alonso, Pexels)}};\end{tikzpicture}}

\begin{document}

% =====================================================================
\maketitle

% ---------------------------------------------------------------------
\begin{frame}{Roteiro}
\begin{columns}[T]
\column{0.5\textwidth}
\begin{enumerate}
  \item[\color{lavanda}\faBed] Por que falar de sono no trabalho
  \item[\color{lavanda}\faBalanceScale] O que a NR-1 passou a exigir
  \item[\color{lavanda}\faSeedling] Do INTEGRA App ao piloto Sono
  \item[\color{lavanda}\faWhatsapp] O canal de conversa no WhatsApp
\end{enumerate}
\column{0.5\textwidth}
\begin{enumerate}
  \item[\color{lavanda}\faChartBar] Painel de gestão e alertas preventivos
  \item[\color{lavanda}\faProjectDiagram] Arquitetura e engenharia de dados
  \item[\color{lavanda}\faUserShield] Tratamento de dados e LGPD
  \item[\color{lavanda}\faFlask] Evidências, mercado e desenho do piloto
\end{enumerate}
\end{columns}
\vfill
\begin{tcolorbox}[colback=nevoa,colframe=nevoa,arc=3pt,left=6pt]
\small\faInfoCircle\ O INTEGRA App foi pensado para cinco eixos de saúde.
Este piloto cobre \textbf{apenas Sono \& Recuperação}. Os demais eixos
ficam para depois, quando o piloto tiver mostrado o que funciona.
\end{tcolorbox}
\end{frame}

% =====================================================================
\section{Por que falar de sono no trabalho}

\begin{frame}{Dormir mal custa caro, e a conta chega ao trabalho}
\begin{columns}[T,onlytextwidth]
\column{0.24\textwidth}
\kpi{coral}{\faDollarSign}{US\$ 411 bi}{perda anual estimada nos EUA
  com sono insuficiente, equivalente a 2,28\,\% do PIB \parencite{hafner2016}}
\column{0.24\textwidth}
\kpi{ambar}{\faHardHat}{1,62$\times$}{risco relativo de acidente de
  trabalho entre pessoas com problemas de sono, que respondem por cerca de 13\,\% dos acidentes \parencite{uehli2014}}
\column{0.24\textwidth}
\kpi{lavanda}{\faBusinessTime}{7,8 dias}{de desempenho perdidos por ano
  por trabalhador com insônia, o que equivale a US\$ 2.280 por pessoa \parencite{kessler2011}}
\column{0.24\textwidth}
\kpi{menta}{\faClock}{$\geq$\,7 h}{de sono por noite é a recomendação
  para adultos \parencite{watson2015,hirshkowitz2015}}
\end{columns}
\bigskip
\small
Esses números não querem dizer que cada pessoa que dorme mal vai render
menos ou se acidentar. Eles mostram uma associação consistente, observada
em muitos estudos, entre sono insuficiente e perdas que a empresa sente:
erros, faltas, presença sem rendimento e acidentes.
\end{frame}

% ---------------------------------------------------------------------
\begin{frame}{Como o sono ruim chega ao trabalho (e o trabalho chega ao sono)}
\centering
\begin{adjustbox}{max width=\textwidth}
\begin{tikzpicture}[node distance=4mm and 9mm]
  % Coluna 1: organização do trabalho
  \node[caixa=noite,text width=3.4cm] (o1) {Turnos noturnos e rodízio};
  \node[caixa=noite,text width=3.4cm,below=of o1] (o2) {Jornadas longas e horas extras};
  \node[caixa=noite,text width=3.4cm,below=of o2] (o3) {Sobrecarga e pressão por metas};
  \node[caixa=noite,text width=3.4cm,below=of o3] (o4) {Demandas fora do horário (mensagens, telas)};
  \node[grupo=noite,fit=(o1)(o4),label={[rot=noite]above:Organização do trabalho}] (g1) {};
  % Coluna 2: sono
  \node[caixa=lavanda,text width=3.0cm,right=14mm of o1] (s1) {Duração};
  \node[caixa=lavanda,text width=3.0cm,below=of s1] (s2) {Qualidade e continuidade};
  \node[caixa=lavanda,text width=3.0cm,below=of s2] (s3) {Regularidade de horários};
  \node[caixa=lavanda,text width=3.0cm,below=of s3] (s4) {Sonolência diurna};
  \node[grupo=lavanda,fit=(s1)(s4),label={[rot=lavanda]above:Saúde do sono}] (g2) {};
  % Coluna 3: desfechos
  \node[caixa=coral,text width=3.3cm,right=14mm of s1] (d1) {Atenção, memória e erros};
  \node[caixa=coral,text width=3.3cm,below=of d1] (d2) {Acidentes e quase-acidentes};
  \node[caixa=coral,text width=3.3cm,below=of d2] (d3) {Absenteísmo e presenteísmo};
  \node[caixa=coral,text width=3.3cm,below=of d3] (d4) {Saúde mental e afastamentos};
  \node[grupo=coral,fit=(d1)(d4),label={[rot=coral]above:Desfechos ocupacionais}] (g3) {};
  % Setas
  \draw[seta=noite] (g1.east) -- (g2.west);
  \draw[seta=lavanda] (g2.east) -- (g3.east -| g3.west);
  \draw[seta=mudo,dashed] (g3.south) to[bend left=18]
     node[below,font=\tiny\color{mudo},text width=6cm,align=center]
     {a relação é bidirecional: estresse e adoecimento também pioram o sono}
     (g1.south);
\end{tikzpicture}
\end{adjustbox}
\fonte{\textcite{linton2015,kecklund2016,uehli2014}. Associações, não
  relações determinísticas.}
\end{frame}

% ---------------------------------------------------------------------
\begin{frame}{E no Brasil?}
\begin{columns}[T,onlytextwidth]
\column{0.58\textwidth}
\small
\begin{itemize}
  \item Em São Paulo, a queixa de acordar durante a noite passou de
    15,8\,\% para 36,5\,\% dos adultos entre 1987 e 2007 \parencite{santossilva2010}.
  \item Num inquérito nacional com 2.635 adultos, 65,5\,\% foram
    classificados como maus dormidores pelo PSQI \parencite{drager2022}.
  \item Afastamentos por transtornos mentais no INSS: 472 mil em 2024 (alta
    de 68\,\%) e 546 mil em 2025 \parencite{barros2025_infomoney,prates2026_dgabc}.
    O sono não é a causa desses números, mas costuma piorar junto com eles.
  \item Trabalho em turnos é rotina em saúde, indústria, logística,
    segurança e serviços. É justamente nesses setores que sono curto e
    acidentes andam juntos \parencite{kecklund2016}.
\end{itemize}
\column{0.38\textwidth}
\centering
\imagem{\linewidth}{4.4cm}{img/trabalhadora_turno_noturno_brasil.jpg}{Profissional de saúde em turno noturno (Foto: Mikhail Nilov, Pexels)}
\end{columns}
\end{frame}

% =====================================================================
\section{O que a NR-1 passou a exigir}

\begin{frame}{A linha do tempo da NR-1}
\centering
\begin{adjustbox}{max width=\textwidth}
\begin{tikzpicture}
  \draw[line width=3pt,color=noite!25,line cap=round] (0,0) -- (15.2,0);
  \marco{0.6}{noite}{\faFile*}{2020}{Nova NR-1 cria o GRO e o PGR}{above}
  \marco{3.6}{lavanda}{\faBrain}{ago. 2024}{Portaria MTE 1.419 inclui fatores de risco psicossociais no GRO}{below}
  \marco{6.6}{menta}{\faAward}{2024}{Lei 14.831 cria o certificado Empresa Promotora da Saúde Mental}{above}
  \marco{9.6}{ambar}{\faChalkboardTeacher}{mai. 2025}{Primeiro ano com caráter educativo e orientativo}{below}
  \marco{12.6}{coral}{\faGavel}{26 mai. 2026}{Exigência plena e fiscalização dos riscos psicossociais}{above}
  \marco{14.9}{noite}{\faMapMarker*}{out. 2026}{Proposta do piloto INTEGRA Sono}{below}
\end{tikzpicture}
\end{adjustbox}
\vfill
\small Hoje a empresa já precisa \textbf{identificar, avaliar, controlar e
registrar} os riscos psicossociais no Programa de Gerenciamento de Riscos,
como faz com ruído ou produtos químicos.
\fonte{\textcite{mte_nr1_2025,mte_portaria1419_2024,mte_portaria765_2025,lei14831_2024,mte_guia_psicossociais_2025}.}
\end{frame}

% ---------------------------------------------------------------------
\begin{frame}{Onde o sono entra no gerenciamento de riscos}
\begin{columns}[c,onlytextwidth]
\column{0.47\textwidth}
\centering
\begin{tikzpicture}
  \foreach \a/\cor/\t [count=\i] in {90/noite/Identificar\\perigos,
      18/lavanda/Avaliar\\riscos, -54/menta/Controlar\\(plano\\de ação),
      -126/ambar/Acompanhar, 162/coral/Registrar\\(inventário)}
    \node[etapa=\cor] (e\i) at (\a:2.15cm) {\t};
  \foreach \i/\j in {1/2,2/3,3/4,4/5,5/1}
    \draw[seta=mudo] (e\i) to[bend left=14] (e\j);
  \node[font=\small\bfseries,color=noite,align=center] at (0,0) {GRO\\PGR};
\end{tikzpicture}
\column{0.5\textwidth}
\small
A NR-1 olha para a \textbf{organização do trabalho}: jornada, turnos,
metas, conflitos, falta de suporte. O objetivo não é diagnosticar pessoas.

\medskip
O sono funciona bem como \textbf{sinal de acompanhamento}. Ele reage rápido
a mudanças de escala, sobrecarga e horas extras, e pode ser medido sem
exames, em poucos segundos, pelo próprio trabalhador.

\medskip
\begin{block}{O que o piloto entrega ao PGR}
\scriptsize Indicadores agregados por setor e turno, histórico das
medidas adotadas e evidência de que elas foram acompanhadas.
\end{block}
\end{columns}
\fonte{\textcite{mte_nr1_2025,mte_guia_psicossociais_2025}.}
\end{frame}

% ---------------------------------------------------------------------
\begin{frame}{Fatores psicossociais que afetam o sono: o que observar e o que fazer}
\scriptsize
\begin{tabularx}{\textwidth}{@{}>{\bfseries}p{3.1cm} L L@{}}
\toprule
Fator de risco (Guia MTE) & Sinal que o piloto acompanha & Exemplos de medida de controle \\
\midrule
Sobrecarga & Queda de recuperação em semanas de pico e relatos de \enquote{cabeça que não desliga} & Teto de horas extras, revisão de metas, redistribuição de demanda \\
Baixo controle e autonomia & Sonolência maior em quem não escolhe escala ou turno & Participação na montagem de escalas e rodízio no sentido horário \\
Falta de suporte & Piora coletiva sem mudança de escala & Escuta da equipe e apoio da liderança \\
Baixas recompensas & Piora associada a ciclos de avaliação e metas & Revisão de critérios e reconhecimento \\
Trabalho remoto ou isolado & Atividade de trabalho perto da hora de dormir & Política de desconexão e envio programado de mensagens \\
\bottomrule
\end{tabularx}
\fonte{Fatores do guia do MTE \parencite{mte_guia_psicossociais_2025}, que traz 13 exemplos e não cita o sono
  diretamente. O vínculo com o sono vem de \textcite{linton2015}: alta demanda, baixo controle e
  pouco apoio predizem distúrbios de sono futuros.}
\end{frame}

% =====================================================================
\section{Do INTEGRA App ao piloto Sono}

\begin{frame}{A lógica do INTEGRA, aplicada ao sono}
\begin{columns}[c,onlytextwidth]
\column{0.52\textwidth}
\centering
\begin{adjustbox}{max width=\linewidth}
\begin{tikzpicture}[passo/.append style={minimum width=1.9cm}]
  % Elipse: eixo x 3.8cm, eixo y 2.3cm; dá espaço aos dois nós da base
  \foreach \t/\cor [count=\i] in {Perceber/noite,Entender/noite!80!lavanda,
      Escolher/lavanda,Experimentar/lavanda!70!menta,Registrar/menta,
      Acompanhar/menta!70!ambar,Ajustar/ambar}
    \node[passo=\cor] (p\i) at ({3.8*cos(90-(\i-1)*360/7)},{2.3*sin(90-(\i-1)*360/7)}) {\t};
  % Setas como arcos elípticos no vão entre nós vizinhos, sentido horário
  \foreach \i in {0,...,6}
    \draw[seta=mudo] ({3.8*cos(90-\i*360/7-17)},{2.3*sin(90-\i*360/7-17)})
      arc[start angle={90-\i*360/7-17},end angle={90-(\i+1)*360/7+17},
          x radius=3.8cm,y radius=2.3cm];
  \node[align=center,font=\scriptsize,text width=2.4cm] at (0,0)
    {{\color{lavanda}\Large\faMoon}\\ \textbf{Sono \&\\Recuperação}};
\end{tikzpicture}
\end{adjustbox}
\column{0.46\textwidth}
\small
O INTEGRA transforma conhecimento em pequenas ações que a pessoa
consegue repetir. No piloto, cada passo vira algo concreto:
\begin{itemize}\scriptsize
  \item \textbf{Perceber}: check-in matinal de 10 segundos.
  \item \textbf{Experimentar}: um experimento de 7 dias escolhido pela
    pessoa (ex.: cafeína só até as 16h).
  \item \textbf{Acompanhar}: tendências pessoais, nunca um \enquote{score}
    de saúde.
  \item \textbf{Ajustar}: manter, adaptar ou abandonar a estratégia, sem
    culpa.
\end{itemize}
\end{columns}
\end{frame}

% ---------------------------------------------------------------------
\begin{frame}{Quatro regras de design que não negociamos}
\begin{columns}[T,onlytextwidth]
\column{0.49\textwidth}
\begin{block}{\faStethoscope\ Monitorar não é diagnosticar}
\scriptsize O robô nunca diz \enquote{você tem insônia}. Ele diz
\enquote{você registrou três manhãs com recuperação abaixo do habitual}
e indica quando vale procurar um profissional.
\end{block}
\begin{block}{\faHourglassHalf\ 20 a 40 segundos por dia}
\scriptsize As respostas são dadas por botões. Se a coleta cansar, a adesão cai na
segunda semana e os dados perdem valor.
\end{block}
\column{0.49\textwidth}
\begin{block}{\faUsers\ Sem ranking e sem culpa}
\scriptsize Nada de \enquote{quem dorme melhor}. Interromper um
experimento não quebra a sequência: experimentos servem para descobrir o
que cabe na rotina.
\end{block}
\begin{block}{\faQuestionCircle\ \enquote{Por que estou recebendo isso?}}
\scriptsize Toda mensagem explica sua origem: um lembrete que a pessoa
configurou ou um experimento que ela escolheu. Ninguém recebe
sugestão sem saber de onde ela veio.
\end{block}
\end{columns}
\end{frame}

% ---------------------------------------------------------------------
\begin{frame}{Por que WhatsApp e não um aplicativo próprio}
\begin{columns}[T,onlytextwidth]
\column{0.5\textwidth}
\textcolor{menta}{\faThumbsUp}\ \textbf{A favor}
\begin{itemize}\footnotesize
  \item Já está no celular de quase todo trabalhador, por isso não exige
    download nem senha nova.
  \item Alcança quem não fica no computador: operação, campo, turno
    noturno, prestadores de serviço.
  \item Botões reduzem a coleta a um toque, no momento em que a pessoa
    vive a experiência, com menos viés de memória \parencite{shiffman2008}.
  \item A conversa é o formato natural para um tom acolhedor.
\end{itemize}
\column{0.5\textwidth}
\textcolor{coral}{\faExclamationTriangle}\ \textbf{Cuidados}
\begin{itemize}\footnotesize
  \item As mensagens passam pela infraestrutura da Meta, o que exige
    minimização de conteúdo e cláusulas de operador
    \parencite{lgpd2018}.
  \item Fora da janela de 24 horas, só podem ser enviados modelos
    aprovados, a quem deu opt-in
    \parencite{meta_servicewindow,meta_optin}.
  \item A política da Meta restringe pedir informação de saúde onde a lei
    exige sistemas reforçados, o que torna necessário um parecer jurídico \parencite{whatsapp_policy}.
  \item O número é pessoal e mistura vida e trabalho, por isso a própria
    pessoa escolhe os horários de envio.
\end{itemize}
\end{columns}
\end{frame}

% =====================================================================
\section{O canal de conversa no WhatsApp}

\begin{frame}{Uma conversa de dez segundos}
\begin{columns}[c,onlytextwidth]
\column{0.4\textwidth}
\centering
\begin{tikzpicture}[scale=0.84,transform shape,
  bot/.style={fill=white,rounded corners=4pt,text width=3.0cm,inner sep=4pt,
              font=\tiny,align=left,blur shadow={shadow blur steps=4,shadow xshift=0.3pt,shadow yshift=-0.3pt}},
  usr/.style={fill=wauser,rounded corners=4pt,inner sep=4pt,font=\tiny,
              align=right,blur shadow={shadow blur steps=4,shadow xshift=0.3pt,shadow yshift=-0.3pt}},
  btn/.style={draw=wa!50,fill=white,rounded corners=3pt,font=\tiny\color{wa},
              minimum width=0.95cm,inner sep=2pt}]
  % Aparelho
  \fill[rounded corners=14pt,black!85] (-0.25,-0.3) rectangle (4.45,7.85);
  \fill[wabg] (0,0) rectangle (4.2,7.0);
  \fill[wa] (0,7.0) rectangle (4.2,7.6);
  \node[anchor=west,font=\tiny\bfseries\color{white}] at (0.15,7.3)
     {\faArrowLeft\quad\faMoon\ INTEGRA Sono};
  \node[anchor=east,font=\tiny\color{white}] at (4.05,7.3) {\faEllipsisV};
  \coordinate (R) at (4.05,0);
  % Mensagens
  \node[bot,anchor=north west] (m1) at (0.15,6.85)
    {Bom dia, Ana! \faSun[regular]\\ Como foi sua recuperação esta noite?};
  \node[btn,below=1mm of m1.south west,anchor=north west] (b1) {Boa};
  \node[btn,right=1mm of b1] (b2) {Razoável};
  \node[btn,right=1mm of b2] (b3) {Difícil};
  \node[usr,anchor=north east] (m2) at ([yshift=-2mm]R |- b1.south)
    {Razoável \textcolor{wa}{\faCheckDouble}};
  \node[bot,anchor=north west] (m3) at ([yshift=-2mm]0.15,0 |- m2.south)
    {Obrigado! De 0 a 10, quanta sonolência você sente agora?};
  \node[usr,anchor=north east] (m4) at ([yshift=-2mm]R |- m3.south)
    {6 \textcolor{wa}{\faCheckDouble}};
  \node[bot,anchor=north west] (m5) at ([yshift=-2mm]0.15,0 |- m4.south)
    {Registrado \textcolor{menta}{\faCheck}\\[2pt]
     Foi a 3ª manhã da semana com recuperação abaixo do seu habitual.
     Quer uma estratégia de 2 min para hoje à noite?};
  \node[btn,below=1mm of m5.south west,anchor=north west,minimum width=1.4cm] (b4) {Quero};
  \node[btn,right=1mm of b4,minimum width=1.4cm] {Agora não};
  \node[font=\tiny\color{mudo},anchor=south] at (2.1,0.05)
    {\faQuestionCircle\ Por que recebo isso?};
\end{tikzpicture}
\column{0.56\textwidth}
\small
\begin{itemize}
  \item São duas perguntas com botão e uma devolutiva. A pessoa termina
    em menos de dez segundos.
  \item A devolutiva compara a pessoa \textbf{com ela mesma}, não com
    colegas nem com uma norma.
  \item O convite para a estratégia é opcional. \enquote{Agora não} é uma
    resposta válida e não gera cobrança.
  \item Linguagem de observação (\enquote{você registrou}), nunca de
    julgamento (\enquote{seu sono está ruim}).
\end{itemize}
\fonte{Mensagens de saúde personalizadas e com frequência ajustável
  funcionam melhor \parencite{head2013}.}
\end{columns}
\end{frame}

% ---------------------------------------------------------------------
\begin{frame}{A jornada ao longo do dia, da semana e do mês}
\centering
\begin{adjustbox}{max width=\textwidth}
\begin{tikzpicture}
  \fill[noite!6,rounded corners=4pt] (-0.3,-0.9) rectangle (15.6,1.25);
  \node[font=\scriptsize\bfseries,color=noite,anchor=west] at (-0.2,1.0) {Diário};
  \node[ev=ambar] (d1) at (1.6,0) {\faSun\ \textbf{07h}\\Check-in de recuperação\\(10 s)};
  \node[ev=menta] (d2) at (5.5,0) {\faCoffee\ \textbf{15h30}\\Lembrete do experimento\\(se ativado)};
  \node[ev=lavanda] (d3) at (9.4,0) {\faMoon\ \textbf{21h45}\\Modo desaceleração:\\1 estratégia de 2 min};
  \node[ev=noite] (d4) at (13.3,0) {\faToggleOff\ \textbf{Depois}\\\enquote{Agora deixe o celular\\descansar também.}};
  \foreach \a/\b in {d1/d2,d2/d3,d3/d4} \draw[seta=mudo] (\a) -- (\b);
  \fill[lavanda!6,rounded corners=4pt] (-0.3,-3.4) rectangle (15.6,-1.2);
  \node[font=\scriptsize\bfseries,color=lavanda,anchor=west] at (-0.2,-1.45) {Semanal e mensal};
  \node[ev=lavanda] (s1) at (1.6,-2.5) {\faFlask\ \textbf{Semana 1}\\Escolha do experimento de 7 dias};
  \node[ev=lavanda] (s2) at (5.5,-2.5) {\faCheckSquare\ \textbf{Dia 7}\\\enquote{Funcionou na\\minha rotina?}};
  \node[ev=menta] (s3) at (9.4,-2.5) {\faChalkboardTeacher\ \textbf{Mensal}\\Palestra INTEGRA Edu\\+ quiz de 3 perguntas};
  \node[ev=coral] (s4) at (13.3,-2.5) {\faClipboardList\ \textbf{Semanas 0, 6, 12}\\Questionários breves\\(ISI, Epworth e PSQI)};
  \foreach \a/\b in {s1/s2,s2/s3,s3/s4} \draw[seta=mudo] (\a) -- (\b);
\end{tikzpicture}
\end{adjustbox}
\vfill
\small Todos os horários são escolhidos pelo participante na adesão e podem
ser alterados a qualquer momento com uma mensagem.
\end{frame}

% ---------------------------------------------------------------------
\begin{frame}{O que medimos e com que instrumentos}
\scriptsize
\begin{tabularx}{\textwidth}{@{}l L c c L@{}}
\toprule
\textbf{Instrumento} & \textbf{O que mede} & \textbf{Itens} & \textbf{Frequência} & \textbf{Leitura no piloto} \\
\midrule
Diário de sono (consenso), versão curta & Qualidade percebida, recuperação, horários & 2 a 4 & Diária & Tendência individual e base dos agregados \\
Índice de Gravidade de Insônia (ISI) & Gravidade dos sintomas de insônia & 7 & Sem.\ 0, 6, 12 & $\geq$\,15 indica insônia clínica moderada, e o estudo brasileiro usou $\geq$\,8 \\
Escala de Sonolência de Epworth & Sonolência diurna em situações do dia a dia & 8 & Sem.\ 0, 6, 12 & $>$\,10 indica sonolência excessiva \\
Índice de Qualidade do Sono de Pittsburgh (PSQI) & Qualidade do sono no último mês & 19 & Sem.\ 0 e 12 & $>$\,5 indica má qualidade \\
\bottomrule
\end{tabularx}
\medskip

\small Usamos versões validadas no Brasil quando existem. Nenhum escore é
devolvido como diagnóstico: valores acima do corte geram apenas um
convite para conversar com um profissional de saúde.
\fonte{Diário: \textcite{carney2012}. ISI: \textcite{bastien2001,castro2011}.
  Epworth: \textcite{johns1991,bertolazi2009}. PSQI: \textcite{buysse1989,bertolazi2011}.}
\end{frame}

% ---------------------------------------------------------------------
\begin{frame}{Como o robô conversa: regras primeiro, linguagem depois}
\begin{columns}[T,onlytextwidth]
\column{0.53\textwidth}
\centering
\begin{adjustbox}{max width=\linewidth}
\begin{tikzpicture}[node distance=5mm and 6mm]
  \node[caixa=noite,text width=2.4cm] (in) {Mensagem recebida};
  \node[diamond,draw=coral,fill=coral!10,aspect=2,inner sep=1pt,
        font=\tiny,align=center,below=of in] (risco) {Sinal de\\crise?};
  \node[caixa=coral,text width=2.6cm,right=8mm of risco] (crise)
     {Protocolo de segurança\\CVV 188 \textperiodcentered\ SAMU 192\\+ equipe de saúde};
  \node[caixa=lavanda,text width=2.4cm,below=of risco] (fsm)
     {Máquina de estados\\(check-in, experimento, questionário)};
  \node[caixa=menta,text width=2.6cm,right=8mm of fsm] (llm)
     {LLM com base curada\\(dúvidas abertas)};
  \node[caixa=noite,text width=2.4cm,below=of fsm] (out) {Resposta revisada\\por filtros de saída};
  \draw[seta] (in) -- (risco);
  \draw[seta=coral] (risco) -- node[rot=coral,above]{sim} (crise);
  \draw[seta] (risco) -- node[rot=grafite,right]{não} (fsm);
  \draw[seta=menta] (fsm) -- node[rot=menta,above]{texto livre} (llm);
  \draw[seta] (fsm) -- (out);
  \draw[seta=menta] (llm) |- (out);
\end{tikzpicture}
\end{adjustbox}
\column{0.44\textwidth}
\scriptsize
\begin{tabularx}{\linewidth}{@{}L L@{}}
\toprule
\textcolor{menta}{\faCheck} \textbf{O robô faz} & \textcolor{coral}{\faTimes} \textbf{O robô não faz} \\
\midrule
Coleta por botões, com fluxo fixo & Diagnóstico ou prescrição \\
Explica com a base \enquote{O que sabemos / O que não podemos concluir / Quando procurar ajuda} & Recomenda remédio para dormir \\
Encaminha para profissional quando há critério & Responde a temas fora do sono \\
Registra a origem de cada sugestão & Repassa a conversa ao gestor \\
\bottomrule
\end{tabularx}
\end{columns}
\fonte{Agentes conversacionais em saúde têm resultados promissores, mas
  heterogêneos e com riscos de segurança quando geram texto livre sem
  controle \parencite{laranjo2018,abdalrazaq2020}.}
\end{frame}

% =====================================================================
\section{Painel de gestão e alertas preventivos}

\begin{frame}{Quem vê o quê}
\scriptsize
\begin{tabularx}{\textwidth}{@{}>{\raggedright\bfseries}p{3.6cm} C C C C@{}}
\toprule
Perfil & Próprios registros & Indicadores agregados (setor, turno) & Dados individuais identificados & Trilhas de auditoria \\
\midrule
Colaborador (no WhatsApp) & \textcolor{menta}{\faCheckCircle} & \nao & só os seus & \nao \\
Gestor / RH & \nao & \textcolor{menta}{\faCheckCircle}\ grupos com $n \geq 10$ & \textcolor{coral}{\faTimesCircle} & \nao \\
Equipe de saúde ocupacional (SESMT, médico do trabalho) & \nao & \textcolor{menta}{\faCheckCircle} & \textcolor{ambar}{\faExclamationCircle}\ só com consentimento e sob sigilo & \nao \\
Encarregado de dados (DPO) & \nao & \textcolor{menta}{\faCheckCircle} & \textcolor{coral}{\faTimesCircle} & \textcolor{menta}{\faCheckCircle} \\
Administrador técnico & \nao & \nao & \textcolor{coral}{\faTimesCircle}\ (dados cifrados) & \textcolor{menta}{\faCheckCircle} \\
\bottomrule
\end{tabularx}
\bigskip

\small A gestão precisa decidir sobre \textbf{escalas, metas e jornadas},
não sobre a noite de uma pessoa específica. Por isso o gestor enxerga
grupos. O acompanhamento individual fica com quem tem dever de sigilo
profissional.
\fonte{\textcite{lgpd2018}, arts.\ 6º, 11 e 46.}
\end{frame}

% ---------------------------------------------------------------------
\begin{frame}{O painel do gestor}
\centering
\begin{adjustbox}{max width=\textwidth,max totalheight=0.8\textheight}
\begin{tikzpicture}[
  tile/.style={fill=white,draw=noite!15,rounded corners=3pt,minimum width=3.0cm,
               minimum height=1.35cm,anchor=north west}]
  % moldura
  \fill[nevoa,rounded corners=5pt] (0,0) rectangle (15.4,-8.3);
  \fill[noite,rounded corners=5pt] (0,0) rectangle (1.3,-8.3);
  \fill[noite] (0.6,0) rectangle (1.3,-8.3);
  \foreach \ic/\y in {\faMoon/-0.6,\faChartLine/-1.6,\faBell/-2.5,
                      \faClipboardList/-3.4,\faUserShield/-4.3,\faCog/-7.6}
    \node[text=white,font=\small] at (0.65,\y) {\ic};
  \node[anchor=west,font=\small\bfseries,color=noite] at (1.6,-0.45)
     {Visão geral \textperiodcentered\ Unidade Norte \textperiodcentered\ últimas 4 semanas};
  \node[anchor=east,font=\tiny,color=mudo] at (15.2,-0.45)
     {\faLock\ dados agregados \textperiodcentered\ grupos com $n\geq10$};
  % KPIs
  \foreach \x/\cor/\v/\l in {1.6/menta/72\%/Adesão (participantes ativos),
      5.05/lavanda/{4,6}/{Check-ins por pessoa por semana},
      8.5/ambar/31\%/Manhãs com recuperação baixa,
      11.95/coral/2/Alertas organizacionais abertos}
  { \node[tile] (t) at (\x,-0.9) {};
    \node[anchor=north west,font=\Large\bfseries,color=\cor] at ([shift={(0.15,-0.1)}]t.north west) {\v};
    \node[anchor=south west,font=\tiny,color=grafite,text width=2.8cm] at ([shift={(0.15,0.08)}]t.south west) {\l}; }
  % Gráfico de barras por turno
  \fill[white,rounded corners=3pt] (1.6,-2.5) rectangle (8.2,-8.0);
  \node[anchor=north west,font=\scriptsize\bfseries] at (1.7,-2.55)
     {Recuperação baixa por turno (\% das manhãs)};
  \begin{axis}[at={(2.3cm,-7.55cm)},anchor=south west,width=6.6cm,height=4.4cm,
     ybar,bar width=7pt,ymin=0,ymax=60,enlarge x limits=0.22,
     symbolic x coords={Manhã,Tarde,Noite,Rodízio},xtick=data,
     tick label style={font=\tiny},ylabel={\%},ylabel style={font=\tiny},
     area legend,
     legend style={font=\tiny,at={(0.5,1.03)},anchor=south,draw=none,
                   legend columns=2,fill=none,/tikz/every even column/.append style={column sep=6pt}},
     axis line style={draw=mudo!50},ymajorgrids,grid style={mudo!20},
     nodes near coords,nodes near coords style={font=\tiny}]
    \addplot[fill=lavanda!50,draw=none] coordinates {(Manhã,22) (Tarde,25) (Noite,41) (Rodízio,36)};
    \addplot[fill=lavanda,draw=none] coordinates {(Manhã,20) (Tarde,24) (Noite,52) (Rodízio,38)};
    \legend{mês anterior,mês atual}
  \end{axis}
  % Mapa de calor setor x semana
  \fill[white,rounded corners=3pt] (8.5,-2.5) rectangle (15.15,-8.0);
  \node[anchor=north west,font=\scriptsize\bfseries] at (8.6,-2.55)
     {Sonolência média (0 a 10) por setor e semana};
  \def\vals{{{3.1,3.4,3.2,3.6},{4.0,4.6,5.3,5.9},{2.8,2.9,-1,-1},{3.5,3.3,3.4,3.1},{4.4,4.9,5.1,5.0}}}
  \foreach \s/\nome [count=\r from 0] in {a/Expedição,b/Manutenção,c/Laboratório,d/Administrativo,e/Portaria}
  { \node[anchor=east,font=\tiny] at (10.75,-3.45-\r*0.8) {\nome};
    \foreach \w in {0,...,3}{
      \pgfmathsetmacro{\v}{\vals[\r][\w]}
      \ifdim\v pt<0pt
        \fill[pattern=north east lines,pattern color=mudo!60] (10.85+\w*1.0,-3.1-\r*0.8) rectangle ++(0.92,-0.7);
        \node[font=\tiny,color=mudo] at (11.31+\w*1.0,-3.45-\r*0.8) {$n<10$};
      \else
        \pgfmathtruncatemacro{\pct}{min(100,max(0,(\v-2.5)*30))}
        \fill[ambar!\pct!menta!35] (10.85+\w*1.0,-3.1-\r*0.8) rectangle ++(0.92,-0.7);
        \node[font=\tiny] at (11.31+\w*1.0,-3.45-\r*0.8) {\v};
      \fi }}
  \foreach \w in {1,...,4}
    \node[font=\tiny,color=mudo] at (10.31+\w*1.0,-7.25) {S\w};
  \node[font=\tiny,color=coral,anchor=west] at (10.85,-7.7) {\faBell\ Manutenção: alta por 3 semanas seguidas};
\end{tikzpicture}
\end{adjustbox}
\fonte{Dados ilustrativos. Células com menos de 10 participantes são
  suprimidas para impedir a identificação de pessoas.}
\end{frame}

% ---------------------------------------------------------------------
\begin{frame}{Alertas preventivos em três níveis}
\scriptsize\renewcommand{\arraystretch}{1.12}
\begin{tabularx}{\textwidth}{@{}l L l L@{}}
\toprule
\textbf{Nível} & \textbf{Gatilho (exemplos)} & \textbf{Quem recebe} & \textbf{Ação esperada} \\
\midrule
\tagc{menta}{1 Pessoal} & 3 manhãs na semana com recuperação abaixo da média pessoal & Só o colaborador & Mensagem educativa e oferta de estratégia \\
\tagc{ambar}{2 Cuidado} & ISI $\geq$\,15, Epworth $>$\,10 ou sonolência alta em função de risco (direção, máquinas) & Colaborador e, se ele consentir, a equipe de saúde & Convite para avaliação e orientação de segurança na tarefa \\
\tagc{coral}{3 Organização} & Setor ou turno com piora persistente por 2 semanas ou mais, $n\geq10$ & Gestor e SESMT & Investigar causas (escala, horas extras, picos) e registrar medida no PGR \\
\tagc{noite}{Crise} & Menção a autolesão ou ideação suicida & Protocolo imediato & CVV 188, SAMU 192 e acionamento da equipe de saúde \\
\bottomrule
\end{tabularx}
\smallskip

\footnotesize O gestor recebe alertas sobre \textbf{grupos}, que pedem decisões
sobre o trabalho. Alertas sobre pessoas seguem pelo caminho do cuidado,
com consentimento, e nunca viram critério de avaliação de desempenho.
\end{frame}

% ---------------------------------------------------------------------
\begin{frame}{Do alerta à decisão estratégica}
\centering
\begin{adjustbox}{max width=\textwidth}
\begin{tikzpicture}[node distance=7mm,
  ]
  \node[pas=lavanda] (a) {\faWaveSquare\\[2pt]\textbf{Sinal}\\Indicador agregado sai da faixa usual};
  \node[pas=lavanda,right=of a] (b) {\faFilter\\[2pt]\textbf{Validação}\\$n\geq10$ e persistência de 2 semanas};
  \node[pas=coral,right=of b] (c) {\faBell\\[2pt]\textbf{Alerta}\\Gestor e SESMT recebem no painel};
  \node[pas=ambar,right=of c] (d) {\faSearch\\[2pt]\textbf{Análise}\\Escalas, horas extras, metas, eventos};
  \node[pas=menta,right=of d] (e) {\faTools\\[2pt]\textbf{Medida}\\Registrada no plano de ação do PGR};
  \foreach \x/\y in {a/b,b/c,c/d,d/e} \draw[seta] (\x) -- (\y);
  \draw[seta=menta,dashed] (e.south) -- ++(0,-0.6) -| node[pos=0.25,below,font=\scriptsize\color{menta}]
     {\faRedo\ Acompanhar: o indicador melhorou nas semanas seguintes?} (a.south);
\end{tikzpicture}
\end{adjustbox}
\vfill
\begin{columns}[T,onlytextwidth]
\column{0.5\textwidth}
\begin{exampleblock}{Exemplo}
\scriptsize Sonolência na Manutenção sobe por três semanas. A análise mostra
plantões extras após uma parada programada. A medida é limitar plantões
consecutivos e acompanhar as quatro semanas seguintes.
\end{exampleblock}
\column{0.46\textwidth}
\small O valor para a gestão está no ciclo fechado de sinal, decisão e
verificação. É esse ciclo que a fiscalização da NR-1 vai querer ver
documentado.
\end{columns}
\end{frame}

% =====================================================================
\section{Arquitetura e engenharia de dados}

\begin{frame}{Arquitetura em camadas}
\centering
\begin{adjustbox}{max width=\textwidth,max totalheight=0.86\textheight}
\begin{tikzpicture}[
  ap/.style={caixa=lavanda,text width=2.05cm,minimum height=1.05cm}]
  % Faixas
  \foreach \y/\h/\cor/\nome in {0.75/1.5/menta/Canal,-0.95/0.9/noite/Borda,
      -2.25/1.7/lavanda/Aplicação,-4.35/1.7/coral/Dados}
  { \fill[\cor!6,rounded corners=4pt] (-2.9,\y) rectangle (14.0,\y-\h);
    \node[rotate=90,font=\scriptsize\bfseries,color=\cor] at (-2.6,\y-\h/2) {\nome}; }
  % Canal
  \node[font=\scriptsize,align=center] (col) at (-1.3,0) {{\color{noite}\Large\faUser}\\Colaborador};
  \node[caixa=menta,text width=1.9cm] (wa) at (1.0,0) {\faWhatsapp\ WhatsApp};
  \node[caixa=menta,text width=2.2cm] (meta) at (3.75,0) {Meta Cloud API\\(operador)};
  \node[caixa=ambar,text width=2.2cm] (front) at (10.0,0) {\faChartBar\ Painel Admin\\(React)};
  \node[font=\scriptsize,align=center] (ger) at (12.7,0) {{\color{ambar}\Large\faUserTie}\\Gestor, SESMT, DPO};
  % Borda
  \node[caixa=noite,minimum width=15.2cm,minimum height=0.6cm,font=\scriptsize] (ngx) at (5.6,-1.4)
    {\faUserShield\ \textbf{Nginx}\quad TLS 1.3 \textperiodcentered\ validação de assinatura do webhook
     \textperiodcentered\ limite de requisições \textperiodcentered\ proxy do login OIDC com MFA};
  % Aplicação
  \node[ap] (gw)  at (-0.6,-3.1) {Gateway\\WhatsApp};
  \node[ap] (fsm) at (1.95,-3.1) {Motor\\conversacional};
  \node[ap] (llm) at (4.5,-3.1) {Serviço de linguagem\\(LLM + filtros)};
  \node[ap] (sch) at (7.05,-3.1) {Agendador de lembretes};
  \node[ap] (rul) at (9.6,-3.1) {Indicadores e alertas};
  \node[ap] (api) at (12.15,-3.1) {API Admin\\(RBAC)};
  % Dados
  \node[dbn=coral]   (id)   at (-0.6,-5.2) {Cofre de identidade (cifrado)};
  \node[dbn=lavanda] (sono) at (2.7,-5.2) {Dados de sono (pseudônimo)};
  \node[dbn=noite]   (fila) at (6.0,-5.2) {Redis\\(fila e agenda)};
  \node[dbn=menta]   (agg)  at (9.3,-5.2) {Analítico\\(agregados)};
  \node[dbn=grafite] (aud)  at (12.15,-5.2) {Auditoria\\(só inclusão)};
  % Ligações
  \draw[seta] (col) -- (wa);
  \draw[seta] (wa) -- (meta);
  \draw[seta] (ger) -- (front);
  \draw[seta=menta] (meta) -- node[rot=menta,right]{HTTPS} (meta |- ngx.north);
  \draw[seta=ambar] (front) -- (front |- ngx.north);
  \draw[seta=menta] (gw |- ngx.south) -- (gw);
  \draw[seta=ambar] (api |- ngx.south) -- (api);
  \draw[seta,<->] (gw) -- (fsm);
  \draw[seta,<->] (fsm) -- (llm);
  \draw[seta] (sch.north) to[bend right=14] (gw.north);
  \draw[seta] (api) -- (rul);
  \draw[seta] (gw) -- (id);
  \draw[seta] (fsm) -- (sono);
  \draw[seta] (sch) -- (fila);
  \draw[seta] (rul) -- (agg);
  \draw[seta] (sono.south) to[out=-25,in=-155] node[rot=menta,below]{ETL agregado} (agg.south);
  \draw[seta] (api) -- (aud);
\end{tikzpicture}
\end{adjustbox}
\fonte{A arquitetura proposta.}
\end{frame}

% ---------------------------------------------------------------------
\begin{frame}{O caminho de uma resposta, passo a passo}
\centering
\begin{adjustbox}{max width=\textwidth,max totalheight=0.82\textheight}
\begin{tikzpicture}[
  vida/.style={draw=mudo!60,dashed},
  lb/.style={font=\tiny,fill=white,inner sep=1pt}]
  \foreach \x/\cor/\t [count=\i] in {0/menta/Colaborador\\(WhatsApp),
      3.2/menta/Meta\\Cloud API, 6.4/noite/Nginx +\\Gateway,
      9.6/lavanda/Motor\\conversacional, 12.8/coral/PostgreSQL}
  { \node[cab=\cor] (h\i) at (\x,0) {\t};
    \draw[vida] (h\i.south) -- (\x,-6.4); }
  \draw[msg] (0,-1.0) -- node[lb,above]{toca \enquote{Razoável}} (3.2,-1.0);
  \draw[msg] (3.2,-1.5) -- node[lb,above]{POST /webhook (assinado)} (6.4,-1.5);
  \draw[msg=coral] (6.4,-2.0) -- node[lb,above,pos=0.7]{troca telefone por pseudônimo (cofre)} (12.8,-2.0);
  \draw[msg] (6.4,-2.5) -- node[lb,above]{evento validado, já sem telefone} (9.6,-2.5);
  \draw[msg=coral] (12.8,-3.0) -- node[lb,above]{grava check-in} (9.6,-3.0);
  \draw[msg] (9.6,-3.5) to[out=-20,in=20,looseness=6] node[lb,right]{próximo estado} (9.6,-4.0);
  \draw[msg] (9.6,-4.6) -- node[lb,above]{texto da próxima pergunta} (6.4,-4.6);
  \draw[msg] (6.4,-5.1) -- node[lb,above]{POST /messages} (3.2,-5.1);
  \draw[msg=menta] (3.2,-5.6) -- node[lb,above]{\enquote{De 0 a 10...}} (0,-5.6);
  \node[font=\tiny,color=mudo,anchor=west,text width=3.2cm] at (13.1,-4.3)
     {Resposta ao webhook em \textless\,1 s,\\ com processamento assíncrono na fila};
\end{tikzpicture}
\end{adjustbox}
\fonte{Fluxo segundo a documentação da WhatsApp Business Platform \parencite{meta_servicewindow}.}
\end{frame}

% ---------------------------------------------------------------------
\begin{frame}{Modelo de dados: identidade separada dos dados de saúde}
\centering
\begin{adjustbox}{max width=\textwidth,max totalheight=0.84\textheight}
\begin{tikzpicture}[
  rel/.style={line width=0.7pt,color=mudo}]
  \node[ent=coral] (pes) at (0,0)
    {\textcolor{white}{\textbf{identidade.pessoa}}
     \nodepart{two}\textbf{id\_pessoa} (PK)\\telefone\_cifrado\\nome\_cifrado\\matricula\_hash};
  \node[ent=coral,below=6mm of pes] (vin) {\textcolor{white}{\textbf{identidade.vinculo}}
     \nodepart{two}id\_pessoa (FK)\\\textbf{pseudo\_id}\\chave em cofre (KMS)};
  \node[ent=lavanda] (par) at (4.1,0)
    {\textcolor{white}{\textbf{sono.participante}}
     \nodepart{two}\textbf{pseudo\_id} (PK)\\setor\_id, turno, unidade\_id\\faixa\_etaria\\status};
  \node[ent=lavanda,below=6mm of par] (con) {\textcolor{white}{\textbf{sono.consentimento}}
     \nodepart{two}pseudo\_id (FK)\\versao\_termo, finalidades\\aceite\_em, revogado\_em};
  \node[ent=lavanda] (chk) at (8.2,0)
    {\textcolor{white}{\textbf{sono.checkin\_diario}}
     \nodepart{two}pseudo\_id (FK), data\\qualidade (1 a 5)\\recuperacao (1 a 3)\\sonolencia (0 a 10)};
  \node[ent=lavanda,below=6mm of chk] (que) {\textcolor{white}{\textbf{sono.questionario}}
     \nodepart{two}pseudo\_id (FK), data\\instrumento (ISI, ESS)\\escore\_total};
  \node[ent=lavanda,below=6mm of que] (exp) {\textcolor{white}{\textbf{sono.experimento}}
     \nodepart{two}pseudo\_id (FK)\\tipo, inicio, fim\\dias\_realizados};
  \node[ent=menta] (agg) at (12.3,0)
    {\textcolor{white}{\textbf{analitico.indicador\_semana}}
     \nodepart{two}setor\_id, turno, semana\\n\_participantes\\media\_sonolencia\\pct\_recuperacao\_baixa\\suprimido (n<10)};
  \node[ent=menta,below=6mm of agg] (ale) {\textcolor{white}{\textbf{analitico.alerta}}
     \nodepart{two}nivel, escopo, gatilho\\papel\_destino, status\\medida\_pgr};
  \node[ent=noite,below=6mm of ale] (aud) {\textcolor{white}{\textbf{auditoria.acesso}}
     \nodepart{two}usuario, papel, recurso\\acao, timestamp, ip};
  \draw[rel] (pes) -- (vin);
  \draw[rel,dashed,color=coral] (vin.east) -- node[rot=coral,above,sloped]{ligação restrita} (par.west);
  \draw[rel] (par) -- (con);
  \draw[rel] (par.east) -- (chk.west);
  \draw[rel] (par.east) -- (que.west);
  \draw[rel] (par.east) -- (exp.west);
  \draw[rel,-{Stealth}] (chk.east) -- node[rot=menta,above]{ETL} (agg.west);
\end{tikzpicture}
\end{adjustbox}
\fonte{Esquemas separados com papéis de banco distintos. Quem acessa
  \texttt{sono} e \texttt{analitico} não tem permissão sobre \texttt{identidade}.}
\end{frame}

% ---------------------------------------------------------------------
\begin{frame}{Pipeline de dados: da mensagem ao indicador}
\centering
\begin{adjustbox}{max width=\textwidth}
\begin{tikzpicture}[node distance=5mm,
  ]
  \node[st=menta] (c) {\faInbox\\[2pt]\textbf{Coleta}\\botões e listas do WhatsApp};
  \node[st=noite,right=of c] (v) {\faCheckDouble\\[2pt]\textbf{Validação}\\faixas, duplicidade, horário};
  \node[st=coral,right=of v] (p) {\faUserSecret\\[2pt]\textbf{Pseudonimização}\\telefone vira pseudo\_id};
  \node[st=lavanda,right=of p] (a) {\faDatabase\\[2pt]\textbf{Armazenamento}\\cifrado em repouso};
  \node[st=menta,right=of a] (g) {\faLayerGroup\\[2pt]\textbf{Agregação}\\semanal por setor e turno};
  \node[st=ambar,right=of g] (s) {\faEyeSlash\\[2pt]\textbf{Supressão}\\grupos com $n<10$ não aparecem};
  \foreach \x/\y in {c/v,v/p,p/a,a/g,g/s} \draw[seta] (\x) -- (\y);
  \node[caixa=noite,below=9mm of g,text width=2.3cm] (pn) {\faChartBar\ Painel};
  \node[caixa=coral,below=9mm of p,text width=4.6cm] (ret)
    {\faTrash*\ Retenção: dados individuais apagados ou anonimizados ao fim
     do piloto, no prazo definido no RIPD};
  \draw[seta] (s.south) |- (pn.east);
  \draw[seta=coral,dashed] (a.south) |- (ret.east);
\end{tikzpicture}
\end{adjustbox}
\vfill
\small
\begin{columns}[T,onlytextwidth]
\column{0.32\textwidth}\textbf{Qualidade}\par\scriptsize Respostas fora de
faixa são rejeitadas no ato, e o robô pede a resposta de novo.
\column{0.32\textwidth}\textbf{Reprodutibilidade}\par\scriptsize Agregações
versionadas em SQL, com testes automatizados.
\column{0.32\textwidth}\textbf{Pesquisa}\par\scriptsize Exportação apenas de
bases anonimizadas, com aprovação ética quando houver publicação.
\end{columns}
\end{frame}

% ---------------------------------------------------------------------
\begin{frame}{Stack tecnológica proposta}
\scriptsize
\begin{tabularx}{\textwidth}{@{}l l L@{}}
\toprule
\textbf{Camada} & \textbf{Tecnologia} & \textbf{Por quê} \\
\midrule
Canal & WhatsApp Business Platform (Cloud API) & API oficial, mensagens interativas, modelos aprovados \\
Borda & Nginx + Let's Encrypt (já em produção) & TLS, limite de requisições, proxy reverso \\
Backend & Python 3.12, FastAPI, SQLModel & Base atual do projeto, que já tem tipagem e testes \\
Filas e agenda & Redis + worker (ex.: Celery ou ARQ) & Lembretes nos horários escolhidos e reenvio automático em caso de falha \\
Linguagem & LLM via API (provedor nacional, ex.: Maritaca Sabiá) & Português nativo, com uso restrito a texto livre e sem dado identificável no prompt \\
Banco & PostgreSQL 16 + pgcrypto, Row-Level Security & Esquemas separados, cifragem de colunas, políticas por papel \\
Painel & React + Vite, biblioteca de gráficos (ex.: Recharts) & SPA leve servida pelo Nginx \\
Identidade & OIDC (ex.: Keycloak) com MFA & Perfis por papel, sessão curta, trilha de acesso \\
Infraestrutura & Docker Compose na VM e backup cifrado diário & Reaproveita o deploy atual e permite migrar para nuvem gerenciada se o uso crescer \\
\bottomrule
\end{tabularx}
\fonte{Escolhas a validar no início do piloto. Nenhuma delas é dependência
  rígida da arquitetura.}
\end{frame}

% =====================================================================
\section{Tratamento de dados e LGPD}

\begin{frame}{Dado de sono é dado pessoal sensível}
\begin{columns}[T,onlytextwidth]
\column{0.42\textwidth}
\small
A LGPD classifica como sensível todo dado \enquote{referente à saúde}
(art.\ 5º, II). Respostas sobre sono, sonolência e insônia entram nessa
categoria. Isso muda a base legal, exige mais segurança e proíbe uso
discriminatório (art.\ 6º, IX).

\medskip
Na relação de trabalho existe assimetria de poder. Um \enquote{aceito}
dado por medo de represália não é consentimento livre. Por isso a
participação é voluntária e quem recusa não sofre nenhuma consequência.
\column{0.55\textwidth}
\scriptsize
\begin{tabularx}{\linewidth}{@{}L L@{}}
\toprule
\textbf{Finalidade} & \textbf{Base legal (a validar com o jurídico)} \\
\midrule
Check-ins e devolutivas ao próprio colaborador & Consentimento específico e destacado (art.\ 11, I) \\
Encaminhamento à equipe de saúde & Consentimento e tutela da saúde por profissional de saúde (art.\ 11, II, f) \\
Indicadores agregados para o PGR & Dados anonimizados (art.\ 12) e cumprimento de obrigação legal da NR-1 (art.\ 11, II, a) \\
Logs de acesso e segurança & Obrigação legal e exercício regular de direitos (art.\ 11, II, a e d) \\
\bottomrule
\end{tabularx}
\end{columns}
\fonte{\textcite{lgpd2018}.}
\end{frame}

% ---------------------------------------------------------------------
\begin{frame}{Privacidade desde a concepção}
\begin{columns}[T,onlytextwidth]
\column{0.32\textwidth}
\begin{block}{\faCompress\ Minimizar}
\scriptsize
\begin{itemize}
  \item Só perguntas sobre sono
  \item Faixa etária em vez de data de nascimento
  \item Mensagens sem nome de doença ou escore
\end{itemize}
\end{block}
\column{0.32\textwidth}
\begin{block}{\faLock\ Proteger}
\scriptsize
\begin{itemize}
  \item TLS em trânsito e AES-256 em repouso
  \item Pseudonimização na entrada
  \item RBAC, MFA e menor privilégio
  \item Auditoria imutável de acessos
\end{itemize}
\end{block}
\column{0.32\textwidth}
\begin{block}{\faBalanceScale\ Governar}
\scriptsize
\begin{itemize}
  \item RIPD antes do início (art.\ 38), porque dado sensível em escala é tratamento de alto risco
  \item Encarregado nomeado (art.\ 41)
  \item Contratos de operador com Meta e provedores
  \item Incidente comunicado à ANPD em 3 dias úteis
\end{itemize}
\end{block}
\end{columns}
\vfill
\begin{tcolorbox}[colback=menta!8,colframe=menta,boxrule=0.6pt,arc=3pt,left=6pt]
\small\faWhatsapp\ \textbf{Direitos do titular dentro da conversa}:
\texttt{MEUS DADOS} envia o que foi registrado, \texttt{PAUSAR} suspende
os lembretes e \texttt{SAIR} revoga o consentimento e dispara a exclusão
(art.\ 18).
\end{tcolorbox}
\fonte{\textcite{lgpd2018,anpd_res15_2024,anpd_ripd,anpd_res2_2022,anpd_guia_seguranca_2021}.}
\end{frame}

% ---------------------------------------------------------------------
\begin{frame}{Adesão e consentimento no WhatsApp}
\centering
\begin{adjustbox}{max width=\textwidth}
\begin{tikzpicture}[node distance=6mm,
  ]
  \node[st2=noite] (a) {\faBullhorn\\\textbf{Convite}\\QR code ou link na palestra, sem cadastro em massa};
  \node[st2=lavanda,right=of a] (b) {\faFileContract\\\textbf{Termo curto}\\em 3 telas, com link para a versão completa};
  \node[st2=lavanda,right=of b] (c) {\faListUl\\\textbf{Escolhas}\\finalidades separadas, cada uma com seu \enquote{sim}};
  \node[st2=menta,right=of c] (d) {\faClock\\\textbf{Preferências}\\horários e tipos de lembrete};
  \node[st2=menta,right=of d] (e) {\faPlay\\\textbf{Início}\\primeiro check-in no dia seguinte};
  \foreach \x/\y in {a/b,b/c,c/d,d/e} \draw[seta] (\x) -- (\y);
  \node[caixa=coral,below=8mm of c,text width=7.6cm] (r)
   {\faUndo\ A qualquer momento, \texttt{SAIR} revoga tudo, sem exigir justificativa
    e sem avisar o gestor};
  \draw[seta=coral,dashed] (r.north -| b.south) -- (b.south);
  \draw[seta=coral,dashed] (r.east) -| (e.south);
\end{tikzpicture}
\end{adjustbox}
\vfill
\small
O consentimento para \textbf{encaminhamento à equipe de saúde} é pedido à
parte, só quando surge a situação. Assim a pessoa decide com o contexto na
mão, e não no primeiro dia.
\end{frame}

% =====================================================================
\section{Evidências, mercado e desenho do piloto}

\begin{frame}{O que a evidência diz sobre intervenções digitais de sono}
\begin{columns}[T,onlytextwidth]
\column{0.6\textwidth}
\footnotesize
\begin{itemize}
  \item A terapia cognitivo-comportamental para insônia em formato digital
    (dCBT-I) reduz o ISI em 5 pontos em média, segundo meta-análise de 33
    ensaios, e o efeito persiste após um ano \parencite{soh2020,zachariae2016}.
  \item Num ensaio com 1.711 pessoas, a dCBT-I também melhorou saúde
    funcional e bem-estar \parencite{espie2019}.
  \item Em trabalhadores de turno, 55\,\% das intervenções no local de
    trabalho aumentaram a duração do sono \parencite{robbins2021,redeker2019}.
  \item Mudar a organização do trabalho, sem falar de sono, rendeu 8 min a
    mais de sono por dia em ensaio controlado \parencite{olson2015}.
  \item Mensagens personalizadas por app reduziram o ISI de 9,2 para 6,8
    em 4 semanas em 116 trabalhadores \parencite{shimamoto2022}, e um
    chatbot de TCC-I melhorou o PSQI \parencite{chiu2025}.
\end{itemize}
\column{0.36\textwidth}
\begin{alertblock}{O que isso \textbf{não} autoriza}
\scriptsize O piloto não é um tratamento de insônia. Ele usa educação e
higiene do sono, e quem precisa é encaminhado para avaliação clínica e
para a TCC-I, recomendada pela diretriz brasileira \parencite{drager2023}.
Essa distinção precisa ficar clara na comunicação.
\end{alertblock}
\end{columns}
\end{frame}

% ---------------------------------------------------------------------
\begin{frame}{Quem já faz algo parecido}
\scriptsize\renewcommand{\arraystretch}{1.12}
\begin{tabularx}{\textwidth}{@{}l L L L@{}}
\toprule
\textbf{Solução} & \textbf{Força principal} & \textbf{O que importamos} & \textbf{Onde o INTEGRA difere} \\
\midrule
Sleepio (Big Health) & dCBT-I testada em ensaios clínicos, com versão prescrita liberada pela FDA & Jornada estruturada e diário de sono & O INTEGRA não é tratamento e foca prevenção e gestão do trabalho \\
Calm / Headspace & Conteúdo curto de sono e relaxamento, com planos para empresas & Microintervenções de 2 minutos & Canal sem app, com dados que alimentam o PGR \\
WHOOP & Recuperação diária calculada com dados de wearable & Ideia de \enquote{prontidão} do dia & Autorrelato validado, sem dispositivo \\
Bearable & Diário de sintomas e correlações pessoais & Check-in rápido, tendência individual & Correlação comunicada como padrão, não causa \\
Wysa & Agente conversacional de TCC com designação \textit{Breakthrough} da FDA & Protocolo de crise e roteiros clínicos & Escopo restrito ao sono \\
\bottomrule
\end{tabularx}
\smallskip

\footnotesize Até onde levantamos, nenhuma dessas soluções conecta o que o
trabalhador relata ao inventário de riscos psicossociais exigido no
Brasil. É nesse ponto que o INTEGRA Sono se posiciona.
\fonte{\textcite{fda_k233577,aasm_sleepiorx_2024,wysa_fda_2022}, além dos sites
  institucionais dos produtos, consultados em out.\ 2026.}
\end{frame}

% ---------------------------------------------------------------------
\begin{frame}{Desenho do piloto}
\begin{columns}[T,onlytextwidth]
\column{0.48\textwidth}
\small
\begin{description}[Duração]\setlength{\itemsep}{2pt}
  \item[Local] 1 unidade, com pelo menos dois turnos
  \item[Público] 150 a 300 colaboradores e prestadores, adesão voluntária
  \item[Duração] 12 semanas de operação
  \item[Desenho] Antes e depois, com medidas nas semanas 0, 6 e 12
  \item[Ética] Protocolo submetido ao CEP, se houver intenção de publicar
\end{description}
\column{0.5\textwidth}
\scriptsize
\begin{tabularx}{\linewidth}{@{}L r@{}}
\toprule
\textbf{Indicador} & \textbf{Meta} \\
\midrule
Adesão entre convidados & $\geq$\,40\,\% \\
Participantes ativos na semana 8 & $\geq$\,60\,\% \\
Check-ins por pessoa por semana & $\geq$\,4 \\
Experimentos de 7 dias concluídos & $\geq$\,1 por pessoa \\
Variação média no ISI (sem.\ 0 $\to$ 12) & redução \\
Alertas de nível 3 com medida registrada & 100\,\% \\
Aceitabilidade (\enquote{recomendaria a um colega}) & $\geq$\,70\,\% \\
\bottomrule
\end{tabularx}
\end{columns}
\fonte{Metas iniciais para discussão. Sem grupo controle, o piloto
  avalia viabilidade e aceitação, não eficácia.}
\end{frame}

% ---------------------------------------------------------------------
\begin{frame}{Cronograma}
\centering
\begin{adjustbox}{max width=\textwidth}
\begin{ganttchart}[
    hgrid,vgrid={*{3}{draw=none},dotted},
    x unit=0.52cm,y unit chart=0.52cm,y unit title=0.55cm,
    title label font=\tiny\bfseries,
    bar label font=\scriptsize,group label font=\scriptsize\bfseries,
    milestone label font=\scriptsize\itshape\color{coral},
    bar/.append style={fill=lavanda!70,draw=none,rounded corners=1pt},
    group/.append style={fill=noite,draw=none},
    milestone/.append style={fill=coral,draw=none},
    bar height=0.55,group peaks height=0.15,group height=0.25]{1}{24}
  \gantttitle{Mês 1}{4}\gantttitle{Mês 2}{4}\gantttitle{Mês 3}{4}
  \gantttitle{Mês 4}{4}\gantttitle{Mês 5}{4}\gantttitle{Mês 6}{4}\\
  \ganttgroup{Governança}{1}{6}\\
  \ganttbar{RIPD, termos, contratos de operador}{1}{4}\\
  \ganttbar{Comitê (RH, SESMT, DPO, CIPA)}{2}{6}\\
  \ganttgroup{Construção}{3}{11}\\
  \ganttbar{Fluxos de conversa e modelos Meta}{3}{7}\\
  \ganttbar{Backend, banco e agregações}{4}{10}\\
  \ganttbar{Painel e alertas}{6}{11}\\
  \ganttbar{Teste com 15 voluntários}{10}{11}\\
  \ganttmilestone{Go/no-go}{11}\\
  \ganttgroup{Operação}{12}{23}\\
  \ganttbar{Palestra de lançamento e adesão}{12}{12}\\
  \ganttbar{12 semanas de piloto}{12}{23}\\
  \ganttbar{Avaliação e relatório}{22}{24}
  \ganttlink{elem6}{elem7}
\end{ganttchart}
\end{adjustbox}
\end{frame}

% ---------------------------------------------------------------------
\begin{frame}{Riscos e como os tratamos}
\scriptsize
\begin{tabularx}{\textwidth}{@{}L L L@{}}
\toprule
\textbf{Risco} & \textbf{Por que importa} & \textbf{Mitigação} \\
\midrule
Medo de vigilância pela chefia & Baixa adesão e respostas \enquote{maquiadas} & Gestor só vê grupos, a comunicação é clara e a CIPA participa do comitê \\
Fadiga de mensagens & Abandono depois da 2ª semana & Horários escolhidos pela pessoa e no máximo 3 mensagens por dia, duas delas opcionais \\
Reidentificação em grupos pequenos & Exposição de dado sensível & Supressão de grupos com $n<10$ e bloqueio de cruzamentos livres de filtros \\
Conselho clínico indevido gerado pelo LLM & Dano ao participante, risco reputacional & Coleta por regras, com LLM restrito à base curada e a filtros de saída \\
Política da Meta sobre dados de saúde & Bloqueio da conta ou do canal & Parecer jurídico prévio, perguntas sem termos clínicos e, como plano B, formulário web próprio acionado pelo WhatsApp \\
Atribuição ao app de efeitos multifatoriais & Conclusões erradas sobre absenteísmo & Avaliação apresentada como estudo de viabilidade, com desfechos tratados como multifatoriais \\
\bottomrule
\end{tabularx}
\end{frame}

% ---------------------------------------------------------------------
\begin{frame}{Para seguir em frente}
\begin{columns}[T,onlytextwidth]
\column{0.58\textwidth}
\small
\textbf{Decisões que dependem da empresa parceira}
\begin{enumerate}\setlength{\itemsep}{0pt}\footnotesize
  \item Unidade e turnos que participarão do piloto
  \item Composição do comitê com RH, SESMT, encarregado de dados e CIPA
  \item Quem da equipe de saúde recebe os encaminhamentos
  \item Aprovação do RIPD e dos termos de consentimento
  \item Conta WhatsApp Business verificada em nome da empresa
\end{enumerate}
\begin{tcolorbox}[colback=noite,colframe=noite,coltext=white,arc=3pt,left=6pt]
\footnotesize Em seis meses, a empresa sai com um indicador de sono por setor e
turno, decisões registradas no PGR e dados para decidir se o INTEGRA
deve crescer para outros eixos.
\end{tcolorbox}
\column{0.38\textwidth}
\centering
\imagem{\linewidth}{4.6cm}{img/equipe_reuniao_bem_estar.jpg}{Equipe em reunião sobre bem-estar no trabalho (Foto: Ivan S, Pexels)}
\end{columns}
\end{frame}

% =====================================================================
\begin{frame}[allowframebreaks]{Referências}
\tiny
\printbibliography[heading=none]
\end{frame}

\end{document}
